\section{Channels of Economic-Model-Quality Tightening:\\
GFM as Quantification of Capability Welfare Economics}
\label{sec:lineage}

Corollary~\ref{cor:direction} states that better economic
models tighten alignment guarantees in the direction of the
operational truth.  As stated, the corollary draws on a single
mechanism (admissible accumulation of revealed-sacrifice
observations), and the argument lives inside a remark wall in
Section~\ref{sec:goodhart_corollary}.  This section develops the
\emph{strong} form of the claim: ``better economic models'' is not
one knob but four, each corresponding to an operational intervention
the deployment can perform on top of \citep[Revealed
Sacrifice]{lasser2026paper8a} and \citep[Need
Sufficiency]{lasser2026paper8b}, each mapping to a specific
component of the proxy-truth gap, and each composing through the
Lipschitz transfer of Theorem~\ref{thm:goodhart} to a
specific alignment-tightening consequence.

The structural lens that organizes the four channels is the
quantification thesis we develop in
Section~\ref{sec:quantification_thesis}: GFM is the
measure-theoretic formalization of the capability-welfare-economics
tradition that runs from Sen 1985 through Stiglitz--Sen--Fitoussi
2009.  The qualitative tradition has been negotiating the same four
moves informally for forty years; what GFM adds is the formal
substrate on which the moves admit Lipschitz analysis,
monotonicity claims, and falsifiable predictions.  The four channels
are exactly the formalized versions of those qualitative moves, and
the Goodhart machinery of Section~\ref{sec:goodhart} is the
alignment-payoff that the quantification makes possible.  Sections
\ref{sec:channel_observation}--\ref{sec:channel_bundles} develop
each channel; Section~\ref{sec:channel_composition} states the
Composition Proposition; Section~\ref{sec:lineage_historical} maps
the lineage onto the channels as historical confirmation that the
moves are real.

\subsection{The quantification thesis}
\label{sec:quantification_thesis}

Welfare economics has been a qualitative discipline.  Sen's
\emph{Commodities and Capabilities}~\citep{sen1985commodities}
introduces the capability set as the appropriate object for
welfare analysis but does not endow it with a measure: the
``capability set'' is a set of functionings the person could
achieve, with comparison rules left to subsequent
work.  Nussbaum's \emph{Women and Human
Development}~\citep{nussbaum2000women} offers a substantive list
of central capabilities, but the list is normative rather than
structural: it does not constrain what counts as a capability or
how capabilities aggregate.  Becker's household-production
mathematics~\citep{becker1965theory,becker1981treatise}
specifies a production function $F(x_1, \dots, x_n)$ for
utility-yielding combinations of inputs, but $F$ is essentially
unconstrained, which makes the framework descriptively powerful and
predictively weak.  The Stiglitz--Sen--Fitoussi commission of
2009~\citep{stiglitz2009report} catalogues at length
\emph{what} GDP misses (non-market production, household labour,
quality-of-life capabilities) but proposes only qualitative
remedies (better surveys, satellite accounts, multi-dimensional
indices) without a unified measure-theoretic framework on which the
remedies compose.

The pattern is recognizable.  Bentham's felicific calculus,
with its qualitative criteria of \emph{intensity, duration,
certainty, propinquity, fecundity, purity, extent}, was a
qualitatively rich descriptive
apparatus, and it took roughly a century-and-a-half for the
expected-utility tradition (Bernoulli, Marshall, Edgeworth, von
Neumann--Morgenstern) to formalize the qualitative claims into a
measure-theoretic apparatus admitting axiomatization, falsification,
and quantitative comparison.  GFM stands to capability welfare
economics in the analogous role:

\begin{quote}
\textbf{The quantification thesis.}  GFM is a measure-theoretic
formalization of the capability-welfare-economics tradition: one
choice of formalization among possible others, not a uniqueness
claim.  The qualitative content of Sen, Nussbaum, Becker, and the
SSF commission is preserved; what is added is the formal substrate
(poset structure, axioms~M1--M6, $\volL$, $\volR$, the
gap decomposition, the residual-class characterization) on which
the qualitative content admits Lipschitz analysis, monotonicity
claims, and falsifiable predictions.  Theorems~\ref{thm:fungibility_collapse}
and~\ref{thm:goodhart} are the alignment-payoff this
formalization buys; the four channels of
Sections~\ref{sec:channel_observation}--\ref{sec:channel_bundles}
are what the formalization makes operationally tractable.
\end{quote}

The thesis is positioning, not a formal theorem of this paper.
What it provides is the interpretive lens through which the
formal results read as connected: the
fungibility-collapse correspondence
(Theorem~\ref{thm:fungibility_collapse}) says where standard utility
theory sits inside the quantified apparatus; the structural
predictions (Section~\ref{sec:structural_predictions}) say what
becomes falsifiable under quantification; the Goodhart theorem
(Theorem~\ref{thm:goodhart}) says how the quantification's
proxy-truth gap controls alignment tightness; and the channels of
Sections~\ref{sec:channel_observation}--\ref{sec:channel_bundles}
say which operational moves drive the gap downward.

\begin{remark}[The qualitative content is preserved, not replaced]
\label{rem:not_replacing}
The quantification thesis does not claim GFM \emph{replaces} the
capability-welfare-economics tradition.  Sen's distinction between
basic and higher capabilities, Nussbaum's normative content,
Becker's household-production reasoning, and SSF's beyond-GDP
agenda all retain their content; what GFM provides is a
measure-theoretic substrate on top of which the existing
content composes.  Specifically, the normative question (which
capabilities should count as central) remains outside the formal
apparatus; the framework provides the scaffold, not the loading.
This is the same relation EU theory has to Bentham's felicific
calculus: EU theory does not tell you \emph{what} to maximize;
it provides the formal conditions under which the maximization is
well-defined.  The present paper inherits the same scope discipline.
\end{remark}

\begin{remark}[What quantification buys, in one inequality]
\label{rem:quantification_buys}
The compact form of what quantification buys is the inequality of
Theorem~\ref{thm:goodhart}:
\[
  | g(T) - g(P) | \leq \mathrm{Lip}(g) \cdot \epsnonres.
\]
The qualitative tradition cannot state this inequality, because
$g$ is not a real-valued function on a measure space and $\epsnonres$
is not a sup-norm of relative deviation between two measures.  The
qualitative tradition can say \emph{the proxy and truth diverge},
\emph{the divergence matters for alignment}, \emph{better measurement
helps}, and the SSF commission has been saying these qualitatively
for nearly two decades.  The inequality is the quantification:
once $g$, $T$, $P$, $\mathrm{Lip}(g)$, and $\epsnonres$ are
formal objects, the qualitative claims become bounds, and the
bounds tighten or loosen with operational interventions in ways
the qualitative tradition could not track.
\end{remark}

The four channels of
Sections~\ref{sec:channel_observation}--\ref{sec:channel_bundles}
are the quantified versions of the qualitative moves the welfare-
economics tradition has been negotiating.  Each section opens with
a one-sentence statement of the qualitative-vs-quantitative pair:
what the tradition implicitly worked with, and what the formal
version is.

\subsection{Channel 1: Observation density}
\label{sec:channel_observation}

\paragraph{Qualitative tradition.}
Welfare economists have long recognized that GDP-style aggregates
miss content; the SSF commission's central recommendation is
``measure what matters,'' which in operational terms means
\emph{collect more data on under-measured dimensions}~\citep{stiglitz2009report}.
The recommendation is qualitative: more is better, with the
direction of improvement specified informally (time-use surveys,
satellite accounts, quality-adjusted indices).  The mathematics
is implicit.

\paragraph{GFM quantification.}
The
\citep[Theorem~2]{lasser2026paper8a} aggregate B-to-C lower bound
$\volRlower$ is what the qualitative recommendation becomes when
the welfare-measurement question is formalized.
\citep[Propositions~2~and~3]{lasser2026paper8a} establish that
$\volRlower$ is non-decreasing in the observation count under
admissible accumulation: each additional admissible event tightens
the lower bound or leaves it unchanged, never weakens it.
$(1 - \betaL)/\betaL$ is then a deployment-readable diagnostic
ceiling on $\epsnonres$ derived from the aggregate lower ratio
$\betaL$, with explicit dependence on the observation count.

\paragraph{Operational intervention.}
The intervention this channel admits: \emph{more admissible
revealed-sacrifice events}.  Concretely, this includes
(i)~increasing the trade-window scope so more events are admissible,
(ii)~adding cross-channel triangulation by composing money-sacrifice
and time-sacrifice ledgers
(\citep[Section~4]{lasser2026paper8a}), (iii)~widening the
deployment's bundle inventory so a larger fraction of the
capability space is in $\delta_{\mathrm{covered}}$ rather than
$\delta_{\mathrm{dormant}}$.

\paragraph{Effect on $\epsgap$.}
The intervention reduces the observation-incomplete share of
$\delta_{\mathrm{covered}}$ in the cell decomposition of
\citep[Definition~8]{lasser2026paper8b}.  Because $\betaL$ is
non-decreasing under admissible accumulation, the diagnostic
ceiling derived from $\betaL$ is non-increasing, and (under
Assumption~\ref{ass:per_cap_admissibility}) $\epsnonres$
admits a non-increasing aggregate bound.  By
Theorem~\ref{thm:goodhart}, the certified upper bound on each
Lipschitz alignment property's deviation from the operational
truth is non-increasing in the same direction.

\paragraph{Where this lands in the lineage.}
The qualitative version of this channel is everything from Sen's
``the capability set should be measured, not just commodities'' to
SSF's catalogue of GDP omissions.  The quantitative version is
$\betaL$ as a non-decreasing diagnostic.  The lineage maps to this
channel through the $\delta_{\mathrm{covered}}$/$\delta_{\mathrm{dormant}}$
distinction: SSF's qualitative ``GDP misses non-market production''
is, in the quantified apparatus, the deployment-specific share of
capability space that has not yet generated admissible trade
events.  More trade events, more cross-channel triangulation, wider
bundle inventory: these are the welfare-economics measurement
moves with explicit alignment-tightening consequences under the
Goodhart machinery.

\subsection{Channel 2: Attestation quality}
\label{sec:channel_attestation}

\paragraph{Qualitative tradition.}
Revealed-preference rigor in the household-production tradition of
Becker~\citep{becker1965theory,becker1981treatise} assumes
voluntary trade events satisfy free-choice and additivity
conditions implicitly: the agent's revealed valuation $\Ui(Y) \geq
\Ui(X)$ is treated as if directly observable from the trade,
without an explicit attestation discipline that classifies events
as voluntary, additively decomposable, and grounded in
benchmarkable counterparts.  In practice, household-production
analyses confront this limitation when revealed-preference data
are mixed with non-voluntary trades (debt-servicing transactions,
forced moves, below-sufficiency need access), and the analyst's
classification of events is qualitative.

\paragraph{GFM quantification.}
\citep[Revealed~Sacrifice]{lasser2026paper8a} states the
attestation requirements as formal assumptions
S0~(valuation-bounded-by-exercise),
S1~(free-choice-with-non-degrading-alternatives), and
S4(a)~(valuation-additivity-on-the-surrendered-side), with the
admissibility condition that events reaching the aggregate satisfy
the named assumptions and the genuine-exchange conditions of
\citep[Definition~2]{lasser2026paper8a}.  These are formal
conditions on each event, not informal stipulations on the data.
\citep[Need~Sufficiency]{lasser2026paper8b} adds the
S1-sufficiency check (Definition~6) that operationalizes the
free-choice condition by requiring at least one non-degrading
alternative in the choice set.  \citep[Remark~18]{lasser2026paper8a}
on external attestation makes explicit that the assumptions are
not zero-knowledge-verifiable under the default commitment
scheme; institutional infrastructure is required to verify them
per-event.

\paragraph{Operational intervention.}
The intervention this channel admits: \emph{sharper attestation of
S0, S1, and S4(a) per admissible event}.  Concretely, this means
institutional infrastructure (trusted-third-party verification,
sybil-resistance protocols, dispute-resolution layers, consent
ledgers) that classifies events with respect to the named
assumptions and admits only those passing the classification.
The four-layer attestation discipline (Layers~A--D) of
\citep[Remark~18]{lasser2026paper8a} catalogues the institutional
preconditions; degradation in any layer reduces the attested
share of admissible events.

\paragraph{Effect on $\epsgap$.}
Sharper attestation reduces the attestation-limited share of
$\delta_{\mathrm{covered}}$: events that would have been excluded
or down-weighted under noisy attestation become admissible
contributions to $\volRlower$ under sharp attestation.  The
mechanism is distinct from Channel~1: Channel~1 increases the
\emph{quantity} of admissible events (more trades pass through);
Channel~2 increases the \emph{certifiable share} of events at any
given quantity.  Both reduce $\epsnonres$ but through different
operational moves, and a deployment's binding-channel question
(``is our gap due to too few events, or to noisy attestation on
the events we have?'') depends on which sub-share of
$\delta_{\mathrm{covered}}$ dominates.

\paragraph{Where this lands in the lineage.}
The qualitative version is Becker's implicit assumption that
revealed-preference data is interpretable as voluntary
exchange, an assumption that fails for debt-driven, coerced, or
below-sufficiency transactions.  The quantitative version is the
S0/S1/S4(a) discipline plus the attestation infrastructure that
verifies it.  The institutional-attestation question is itself an
open one: adversarial degradation of attestation under
optimization pressure (Open~Question~8 of
Section~\ref{sec:discussion}) is the natural sequel to
Conjecture~\ref{conj:optimization_pressure}.

\subsection{Channel 3: Individuation discipline}
\label{sec:channel_individuation}

\paragraph{Qualitative tradition.}
Sen's capability approach treats the capability set as a structural
object whose membership conditions are partly normative and partly
empirical, and the question ``which capabilities count?'' is
addressed in Sen 1985 through ad hoc choices about granularity
and aggregation~\citep{sen1985commodities,sen1979equality}.
Nussbaum's central-capabilities list~\citep{nussbaum2000women}
provides a normative answer (ten capabilities considered central
to a fully human life), but the list is descriptive of important
capabilities rather than structural about \emph{individuation}, i.e., about
when two distinct activities should be treated as one capability or
two.  In welfare-economics practice, the same individuation
ambiguity appears in the choice of consumption basket categories,
the time-use survey's activity taxonomy, and the satellite
accounts' sectoral decomposition.

\paragraph{GFM quantification.}
\citep[Need~Sufficiency]{lasser2026paper8b}, Definition~12,
formalizes the residual class $\ResS$ as the structural set of
capabilities the sacrifice channel cannot in principle reach
(capabilities exercised without any material or
time-displacement footprint).  The complement $P \setminus \ResS$
is the channel-reachable subset on which the framework's
within-channel disagreement gap $\epsnonres$
(Definition~\ref{def:eps_nonres}) is defined.  Whether a given
capability lies in $\ResS$ depends on individuation choices: how
the deployment carves capability space into distinct units, what
counts as a meaningful sacrifice trace, what attestation
infrastructure exists for previously-unsupported activities.
Definition~\ref{def:eps_floor}'s residual-class floor $\epsfloor$
is the share of total operational truth that lives in
$\ResS$ under the deployment's chosen individuation.

\paragraph{Operational intervention.}
The intervention this channel admits: \emph{re-individuation that
moves capabilities from $\ResS$ into the channel-reachable
subset}.  Concretely, this means treating activities that
previously had no recorded sacrifice trace as distinct capabilities
with imputed pricing or imputed time cost: household labour gets
imputed-time pricing via time-use surveys; volunteer activity gets
imputed-wage pricing; previously-uncounted ecosystem services get
satellite-account valuation.  Each re-individuation move pulls a
capability from $\ResS$ into $\delta_{\mathrm{covered}}$ or
$\delta_{\mathrm{dormant}}$ at the cost of admitting an imputation
methodology that the deployment must defend.

\paragraph{Effect on $\epsgap$.}
Re-individuation reduces $\epsfloor$ directly: capabilities that
leave $\ResS$ stop contributing to the floor and start contributing
(via Channels~1, 2, or 4) to the within-channel component
$\epsnonres$.  The total alignment slack $\epsgap = \epsnonres +
\lambda \cdot \epsfloor$ may not reduce monotonically under
re-individuation, because moving a capability from $\ResS$ adds
a new contribution to $\epsnonres$ (the imputation's measurement
slack) at the same time as it subtracts from $\epsfloor$.
Whether the net move tightens depends on whether the imputation's
contribution to $\epsnonres$ is smaller than $\lambda$ times the
floor reduction.  This is the deployment's individuation calculus,
and it is exactly the calculus welfare economics has been
performing implicitly when adopting time-use surveys, satellite
accounts, and imputed-rent calculations.

\begin{remark}[The floor is the binding constraint, eventually]
\label{rem:floor_binds}
For sufficiently mature deployments (those with dense observation
infrastructure (Channel~1), mature attestation (Channel~2), and
fine bundle decomposition (Channel~4)), the within-channel
disagreement $\epsnonres$ may approach the diagnostic-ceiling
floor at which the $\delta_{\mathrm{covered}}$ contribution is
exhausted.  Beyond that point, further tightening of $\epsgap$
requires either (a)~moving capabilities out of $\ResS$ via
re-individuation (treating household labour as a distinct
capability with imputed-time pricing, etc.), or (b)~accepting
the deployment's $\lambda \cdot \epsfloor$ as the operational
floor.  Move~(a) is the GFM-side analogue of welfare economics'
imputation methods~\citep{stiglitz2009report,aguiar2007measuring};
move~(b) is the governance question of which capabilities a
deployment treats as ``no sacrifice-observable path in
principle.''  The floor binds when Channels~1, 2, and 4 saturate;
Channel~3 is what the deployment uses to keep tightening past the
saturation point.
\end{remark}

\paragraph{Where this lands in the lineage.}
The qualitative version is everything from Sen's ad hoc capability
granularity to Nussbaum's central-capabilities list to SSF's
satellite-account programme.  The quantitative version is $\ResS$
as a definable structural class with a deployment-specific
individuation parameter.  Whether GFM's apparatus delivers
sharper imputation tools than the standard welfare-economics
methods is Open~Question~7, a question shared with welfare
economics rather than unique to the alignment programme.

\subsection{Channel 4: Bundle decomposition}
\label{sec:channel_bundles}

\paragraph{Qualitative tradition.}
Becker's household production model~\citep{becker1965theory,becker1981treatise}
specifies utility as $u = u(z_1, \dots, z_m)$ where $z_j =
F_j(x_{j,1}, \dots, x_{j,n_j})$ are produced commodities derived
from inputs.  $F_j$ is essentially unconstrained: the modeller
specifies the production function based on the application, and
the resulting utility maximization over $x$ is descriptively
flexible but predictively soft because $F_j$ admits any
specification.  The empirical content of the framework is
correspondingly limited: most observed bundle behaviour can be
rationalized by the right $F_j$, leaving the framework
descriptively powerful but unfalsifiable in the strong sense.

\paragraph{GFM quantification.}
\citep[Definition~4]{lasser2026paper8a} defines the event-local
bundle decomposition: for each admissible event with bundle $Y_n =
\{y_{n,1}, \dots, y_{n,k_n}\}$, the per-capability coefficients
$\alpha_{n,c}$ partition the bundle's contribution to
$\Delta\volL(X_n)$ across capability components.  The
coefficients are constrained by the poset's structural
connectivity (M6 admits cooperative excess only when the
capabilities are poset-coherent, not arbitrarily so) and by the
within-bundle poset-independence requirement of
\citep[Theorem~2 Part~B]{lasser2026paper8a} when per-capability
aggregation is wanted.  The constraint is what makes Becker's
flexible $F_j$ predictively sharp under quantification:
cooperative capabilities have constrained structure, and the
constraint generates Predictions~\ref{pred:bundle_completion}
and~\ref{pred:coop_premium} of
Section~\ref{sec:structural_predictions} as falsifiable claims.

\paragraph{Operational intervention.}
The intervention this channel admits: \emph{per-component bundle
attribution at the protocol level}.  Concretely, this means
extending the default commitment scheme to expose per-capability
contributions for events whose bundles have non-trivial
cooperative structure.  The
\citep[Open~Question~2]{lasser2026paper8a} flags this as the
bundle-decomposition extension; \citep[Remark~5]{lasser2026paper8a}
notes the bundle-for-bundle generalization as related future work.
Per-event subset-membership openings (the Part-B aggregation
mode of \citep[Theorem~2]{lasser2026paper8a}) and
hedonic-disaggregation panels are the deployment instruments.

\paragraph{Effect on $\epsgap$.}
Bundle decomposition reduces the intra-bundle attribution
mismatch share of $\delta_{\mathrm{covered}}$.  When bundles are
treated as monolithic, per-capability deviations average out
within the bundle, and $\epsnonres$ inherits an aggregate-level
slack that is not separable into per-capability components.
With bundle decomposition, per-capability lower-ratio assumptions
become testable, and $\epsnonres$ admits a per-capability
decomposition that the aggregate-only $\betaL$ cannot deliver.
This is what makes Theorem~\ref{thm:goodhart}'s Lipschitz
transfer cleanly per-property: each alignment property cares about
specific capability components, and the per-component decomposition
lets the Lipschitz constants land on the appropriate components.

\paragraph{Where this lands in the lineage.}
The qualitative version is Becker's $F_j$ as an unconstrained
production function.  The quantitative version is M6's
poset-coherent cooperative excess, with $\alpha_{n,c}$ as the
per-component attribution.  The Becker-to-GFM move is the
constraint-imposition step: by requiring cooperative capabilities
to be poset-coherent, the framework converts Becker's
descriptively-flexible mathematics into a structurally-falsifiable
one.

\subsection{Composition: alignment properties inherit channel sensitivities}
\label{sec:channel_composition}

The four channels are not independent.  Channel~4 (bundle
decomposition) requires sufficient observation density (Channel~1)
and adequate attestation (Channel~2) to be meaningful per-event;
Channel~3 (individuation) is upstream of all three within-channel
channels because it determines which capabilities are in
$\ResS$ versus the channel-reachable subset where the within-channel
gap is even defined.  But under admissible composition, the four
channels each contribute a separate reduction to a separate
component of $\epsgap$; the unclipped aggregate alignment-slack
reduction is the sum of channel contributions, capped at the
pre-intervention slack.

The proposition operates on the \emph{aggregate diagnostic
ceiling} on $\epsnonres$ rather than the sup-norm directly.
Recall from Remark~\ref{rem:agg_vs_sup} that
$(1 - \betaL)/\betaL$ is a deployment-readable diagnostic
ceiling on an aggregate version of $\epsnonres$, additively
decomposable across the cells of
\citep[Definition~8]{lasser2026paper8b}.  Denote this aggregate
diagnostic ceiling by $\hat{\epsilon}^{\mathrm{nonres}}$.  The
sup-norm $\epsnonres$ is bounded above by
$\hat{\epsilon}^{\mathrm{nonres}}$ only under an explicit
admissibility assumption that names how the deployment's
per-capability lower-ratios and channel sub-share allocations
relate to the aggregate.

\begin{assumption}[Per-capability admissibility]
\label{ass:per_cap_admissibility}
A deployment of the framework is \emph{per-capability admissible}
when both of the following hold:
\begin{enumerate}
\item[(i)] \textbf{Per-capability ceiling allocation.}  For every
  $c \in P^{\mathrm{act}}$, the relative deviation
  $|P(c) - T(c)|/T(c)$ is bounded above by a deployment-assigned
  non-negative ceiling component $\hat{\epsilon}_c \geq 0$, and
  the assignment satisfies
  $\sum_{c \in P^{\mathrm{act}}} \hat{\epsilon}_c \leq
   \hat{\epsilon}^{\mathrm{nonres}}$.  Equivalently, the
  aggregate ceiling allocates a non-negative budget across
  per-capability ceilings whose sum bounds the total.
\item[(ii)] \textbf{Disjoint channel sub-shares.}  The channel
  sub-share allocations
  $(\Delta_1, \Delta_2, \Delta_4, \Delta_3^{\mathrm{floor}},
   \Delta_3^{\mathrm{new}})$ act on mutually non-overlapping
  subsets of the per-capability ceiling components in (i): the
  sub-shares have disjoint support across the
  $\delta_{\mathrm{covered}}$, $\delta_{\mathrm{dormant}}$, and
  $\delta_{\mathrm{residual}}$ cells of
  \citep[Definition~8]{lasser2026paper8b}, and within
  $\delta_{\mathrm{covered}}$ the attestation-limited and
  intra-bundle-mismatch portions are themselves disjoint.
\end{enumerate}
Under per-capability admissibility, $\epsnonres \leq
\hat{\epsilon}^{\mathrm{nonres}}$ and the additive
decomposition of channel-wise reductions is well-defined.
\end{assumption}

\begin{remark}[Per-capability admissibility is a deployment
condition, not a structural feature]
\label{rem:per_cap_admissibility_deployment}
Assumption~\ref{ass:per_cap_admissibility} is not a theorem of
the apparatus.  Part~(i) requires the deployment to have
sufficient per-capability ledger granularity to assign
per-capability ceilings; in the public-ledger-only regime where
only event-level aggregates are exposed, the assignment is a
coarse approximation and the sup-norm bound is correspondingly
loose.  Part~(ii) is structurally guaranteed across cells of the
$\delta$-decomposition by
\citep[Definition~8]{lasser2026paper8b}, but the within-cell
disjointness (attestation-limited versus intra-bundle-mismatch
portions of $\delta_{\mathrm{covered}}$) requires the
deployment's classification to separate the two operational
regimes.  When per-capability admissibility fails, the
proposition's bound on $\hat{\epsilon}^{\mathrm{nonres}}$ remains
valid as an aggregate diagnostic but does not bound the sup-norm
directly, and the per-channel attribution becomes a
deployment-readable diagnostic rather than an operationally
tightening claim.
\end{remark}

\begin{proposition}[Channel composition for alignment-tightening]
\label{prop:channel_composition}
Let $g$ be an alignment property satisfying the Lipschitz
hypothesis of Theorem~\ref{thm:goodhart} on
$P^{\mathrm{act}}$, with property-specific Lipschitz constant
$\mathrm{Lip}(g)$.  Assume per-capability admissibility
(Assumption~\ref{ass:per_cap_admissibility}) so that
$\epsnonres \leq \hat{\epsilon}^{\mathrm{nonres}}$.  Suppose a
deployment performs admissible interventions on Channels~1--4
yielding component-wise reductions $\Delta_1, \Delta_2, \Delta_4
\geq 0$ in the corresponding non-overlapping sub-shares of
$\hat{\epsilon}^{\mathrm{nonres}}$ (drawn from
$\delta_{\mathrm{dormant}}$, the attestation-limited portion of
$\delta_{\mathrm{covered}}$, and the intra-bundle mismatch
portion of $\delta_{\mathrm{covered}}$ respectively), with each
$\Delta_i$ bounded above by its allocated component share so
that the post-intervention sub-share remains non-negative.
Channel~3 yields two coupled component-wise quantities:
$0 \leq \Delta_3^{\mathrm{floor}} \leq \epsfloor$, the mass
moved out of $\ResS$ by re-individuation, and
$\Delta_3^{\mathrm{new}} \geq 0$, the new aggregate-ceiling
contribution from the newly-individuated capabilities.  Channel~3 is \emph{admissible} when its
$\lambda$-weighted net is non-negative,
\[
  \Delta_3^{\lambda}
  \;\equiv\; \lambda \cdot \Delta_3^{\mathrm{floor}}
              - \Delta_3^{\mathrm{new}}
  \;\geq\; 0,
\]
which is the condition under which re-individuation tightens the
alignment slack rather than loosening it (cf.\
Remark~\ref{rem:floor_binds} and
Remark~\ref{rem:channel3_net}).  Then the post-intervention
alignment slack admits the bound
\begin{align*}
  \bigl| g(T) - g(P) \bigr|_{\mathrm{post}}
  \leq\;\;& \mathrm{Lip}(g) \cdot \\
  & \Bigl(
    \bigl(\hat{\epsilon}^{\mathrm{nonres}}
       + \Delta_3^{\mathrm{new}}
       - (\Delta_1 + \Delta_2 + \Delta_4)\bigr) \\
  & \;\;{}+ \lambda \cdot (\epsfloor - \Delta_3^{\mathrm{floor}})
  \Bigr)^+,
\end{align*}
where $(\cdot)^+$ denotes the positive part.  The
\emph{unclipped} pre-to-post slack decreases by exactly
$(\Delta_1 + \Delta_2 + \Delta_4) + \Delta_3^{\lambda}$;
the clipped bound's reduction is therefore
\[
  \mathrm{Lip}(g) \cdot \min\!\Bigl(
    \hat{\epsilon}^{\mathrm{nonres}}
      + \lambda \cdot \epsfloor,\;\;
    (\Delta_1 + \Delta_2 + \Delta_4) + \Delta_3^{\lambda}
  \Bigr),
\]
i.e., capped at the pre-intervention slack when the
sum of channel reductions overshoots.  In particular the bound
is non-increasing whenever
$\Delta_1, \Delta_2, \Delta_4 \geq 0$ and
$\Delta_3^{\lambda} \geq 0$.  The composition is exact under
the \emph{sequential-intervention regime} (deployments that
fix one channel before performing the next, so that
Assumption~\ref{ass:per_cap_admissibility}'s disjointness
condition holds across each per-channel step), and approximate
otherwise, with cross-channel correction terms appearing under
simultaneous co-evolution of the channels (see
Remark~\ref{rem:ch34_interaction}).
\end{proposition}

\begin{proof}[Sketch]
Theorem~\ref{thm:goodhart} applied at the post-intervention
configuration gives a within-channel bound
$\mathrm{Lip}(g)$ times the post-intervention sup-norm.  Under
per-capability admissibility, that sup-norm is bounded by the
post-intervention aggregate ceiling, which decomposes additively
as
\[
  \hat{\epsilon}^{\mathrm{nonres}}
   + \Delta_3^{\mathrm{new}}
   - (\Delta_1 + \Delta_2 + \Delta_4),
\]
where the additive decomposition follows from the
non-overlapping sub-share allocation across the cells of
\citep[Definition~8]{lasser2026paper8b}.  The floor side
contributes $\lambda \cdot (\epsfloor - \Delta_3^{\mathrm{floor}})$
under the same per-cell allocation.  Combining and applying the
positive-part operator (which handles the case where the sum of
reductions exceeds the pre-intervention components) gives the
stated bound.
\end{proof}

\begin{remark}[Channel~3's net contribution to the bound]
\label{rem:channel3_net}
The proposition treats Channel~3's contribution as the
$\lambda$-weighted net $\Delta_3^{\lambda} = \lambda \cdot
\Delta_3^{\mathrm{floor}} - \Delta_3^{\mathrm{new}}$ rather than
a pure floor reduction.  This is structurally important:
re-individuation moves a capability $c$ from $\ResS$ into the
channel-reachable subset, which does two things simultaneously.
It reduces $\epsfloor$ by $\vol_R(c)/\vol_R(P)$ (the floor mass
moves out), and it introduces a new contribution to
$\hat{\epsilon}^{\mathrm{nonres}}$ equal to the
aggregate-ceiling share of the newly-individuated capability's
within-channel proxy-truth disagreement.  The two contributions
sit on different sides of the
$\epsgap = \epsnonres + \lambda \cdot \epsfloor$ split: the
floor reduction enters with weight $\lambda$, while the
within-channel addition enters with weight one.  Channel~3 is
non-tightening (the intervention loosens the alignment slack
rather than tightening it) precisely when
$\Delta_3^{\mathrm{new}} > \lambda \cdot
\Delta_3^{\mathrm{floor}}$, equivalently when
$\Delta_3^{\lambda} < 0$.  This is a deployment-specific
calculus (the imputation calculus of
Remark~\ref{rem:floor_binds}), and not a structural feature of
the apparatus.  In the limit $\Delta_3^{\mathrm{new}} = 0$
(the newly-individuated capability has zero within-channel
disagreement on day-one), $\Delta_3^{\lambda} = \lambda \cdot
\Delta_3^{\mathrm{floor}}$ and the proposition reduces to the
pure floor-reduction framing.
\end{remark}


\begin{remark}[Two distinct sources of tightening, formalized]
\label{rem:two_sources}
Proposition~\ref{prop:channel_composition} formalizes a distinction
that Corollary~\ref{cor:direction} conflated.  Channels~1 and~2
deliver \emph{observation-density} and \emph{attestation-quality}
tightening on existing within-channel cells; Channel~3 delivers
\emph{individuation} tightening by moving capabilities between
$\ResS$ and the within-channel subset; Channel~4 delivers
\emph{bundle-decomposition} tightening by sharpening per-capability
attribution.  These are operationally distinct interventions.  When
a deployment reports a tightening, the channel-attribution tells
the deployer which intervention is operative, which is the
information needed to plan the next intervention.  The
five-cell $\delta$-decomposition of
\citep[Definition~8]{lasser2026paper8b} is the operational
instrument that supports the attribution: a deployment-readable
diagnostic that points the deployer toward the binding channel.
\end{remark}

\begin{remark}[Channel sensitivities are property-specific]
\label{rem:channel_property_specific}
Proposition~\ref{prop:channel_composition} states the channel
contributions as separate $\Delta_i$, but the
\emph{property-specific} sensitivity of each alignment property to
each channel is not bounded at this paper's level of analysis.
The anti-monopolar threshold (\citep[Proposition~6]{lasser2026horizon})
cares about substrate-volume distributions and is more sensitive
to Channel~3 (individuation across substrates) than to Channel~4
(within-bundle attribution).  The convergence rate
(\citep[Theorem~2]{lasser2026wamura}) cares about per-capability
exercise dynamics and is more sensitive to Channels~1 and~4 than
to Channel~3.  The phase-boundary location is sensitive to all
four through the Lyapunov function's smoothness profile.
Per-property channel-sensitivity coefficients are perturbative
roadmap items (Section~\ref{sec:goodhart_examples}); the
proposition states the channel-wise composition without
committing to per-property weights.
\end{remark}

\begin{remark}[Channels 3 and 4 are not independent in practice]
\label{rem:ch34_interaction}
Assumption~\ref{ass:per_cap_admissibility} part~(ii) requires the
channel sub-share allocations to act on disjoint cells of
\citep[Definition~8]{lasser2026paper8b}'s decomposition.
Across-cell disjointness (Channels~1, 2 acting on
$\delta_{\mathrm{covered}}$ vs $\delta_{\mathrm{dormant}}$;
Channel~3 acting on $\delta_{\mathrm{residual}}$) is structurally
clean.  But Channels~3 and~4 interact in deployment practice: a
re-individuation move (Channel~3) changes the granularity at
which bundles can be decomposed, which shifts what
$\Delta_4$ is allowed to act on.  Conversely, a sharper bundle
decomposition (Channel~4) can suggest re-individuation choices
that would not be visible without the per-component attribution.
The disjointness condition holds when the deployment fixes its
individuation discipline before performing bundle decomposition
(or vice versa) and treats subsequent moves as separate
interventions; under simultaneous co-evolution of the
individuation scheme and the bundle decomposition,
cross-channel interaction terms appear and the additive
composition of Proposition~\ref{prop:channel_composition} acquires
correction terms.  Characterizing those correction terms, and
the deployment conditions under which they vanish or remain
small, is open (a refinement of
Open~Question~\ref{oq:rate}'s rate question, with the
non-disjointness as a separate sensitivity dimension).  In the
worked examples of Section~\ref{sec:worked_examples} we treat
the channels as performed sequentially, which is the
disjointness regime; deployments that co-evolve Channels~3
and~4 will need the correction-term refinement.
\end{remark}

\subsection{The capability-economics lineage as historical confirmation}
\label{sec:lineage_historical}

The four channels have been catalogued, formalized, and tied to
specific cells of the gap decomposition.  The remaining task is
historical: showing that the channels are not GFM-internal
inventions but the formalized versions of moves the
capability-welfare-economics tradition has been negotiating
qualitatively for forty years.  This subsection traces the
correspondences.

\subsubsection{Sen 1985: the capability set, formalized}

Sen's \emph{Commodities and
Capabilities}~\citep{sen1985commodities} introduces three
structural objects: the \emph{capability set} (the set of
functioning vectors a person could achieve), the \emph{functioning
vector} (what the person actually does or is), and the
\emph{capability--functioning gap} (the difference between
freedom-to-do and actual-doing).  None is endowed with a measure;
the comparison rules between capability sets are addressed
informally and partially.

GFM provides the measure-theoretic counterparts:
\begin{itemize}
\item Sen's \emph{capability set} $\to$ the poset $P$ with measure
  $\volL$ satisfying axioms~M1--M6
  (\citep[Proposition~1]{lasser2026poset}).
\item Sen's \emph{functioning vector} $\to$ the exercised
  sub-poset $P^{\mathrm{ex}}$ with the window-active realized
  exercise volume $\volRwin$
  (\citep[Definition~3]{lasser2026paper8a}).
\item Sen's \emph{capability--functioning gap} $\to$ the B-to-C
  gap with the five-cell $\delta$-decomposition
  (\citep[Definition~8]{lasser2026paper8b}).
\end{itemize}
The functionings--capabilities distinction is recovered through
measurement: $\volL$ measures the option set; $\volR$ measures the
exercised subset; the polarity boundary
(\citep[Definition~3]{lasser2026paper8b}) plus the
S1-attestation layer (\citep[Remark~18]{lasser2026paper8a})
distinguish ``chooses not to exercise'' from ``cannot exercise.''
Sen's freedom dimension is preserved, formally rather than
ontologically: above the polarity boundary, sacrifice reveals
preference; below, sacrifice reveals cost-of-access.

This is the home of all four channels: Channel~1 (observation
density on the capability set), Channel~2 (attestation
infrastructure for the freedom dimension), Channel~3
(individuation of what counts as a capability), and Channel~4
(bundle decomposition within the capability set) all live on the
poset Sen named without a measure.  Sen's later work,
particularly \emph{The Idea of Justice}~\citep{sen2009idea}, is
congenial to the GFM stance on operational truth: Sen's shift
toward a \emph{comparative} rather than transcendental approach
to justice (emphasizing operational comparison of options over
identification of an ideal) parallels GFM's commitment to
operational $T = \volRwin$ rather than metaphysical inner
preference (Definition~\ref{def:operational_T}).  For current
state-of-field surveys of the qualitative tradition,
see~\citep{robeyns2017wellbeing} and~\citep{chiappero2020cambridge}.

\subsubsection{Nussbaum 2000: basic vs higher capabilities, structurally separated}

Nussbaum's \emph{Women and Human
Development}~\citep{nussbaum2000women} adds the basic-versus-higher
capability distinction with substantive normative content.  Basic
capabilities are minimum thresholds for a viable life; higher
capabilities are freedoms the agent can exercise above the
thresholds.  The distinction is qualitative, drawn from
philosophical anthropology rather than structural analysis.

\citep[Need~Sufficiency]{lasser2026paper8b} structurally
separates the two regimes:
\begin{itemize}
\item Below the polarity boundary
  (\citep[Definition~3]{lasser2026paper8b}): need-side, where
  retention is degrading and sacrifice reveals cost-of-access
  (the need-access cost \citep[Definition~9]{lasser2026paper8b}).
\item Above the polarity boundary: want-side, where the agent has
  non-degrading alternatives and sacrifice reveals preference
  (the standard \citep[Revealed~Sacrifice]{lasser2026paper8a}
  channel).
\end{itemize}
The structural separation is what GFM adds beyond Nussbaum:
not just \emph{the basic capability is met} but \emph{the way it
is met does not corrupt the downstream capability structure}
(\citep[Definition~2(b)]{lasser2026paper8b}, the downstream-safety
condition).  Channel~2 (attestation) is most directly relevant to
this distinction: the S1-sufficiency check is what verifies the
above-versus-below polarity classification.  The normative
loading (\emph{which} capabilities are basic, in Nussbaum's
sense) remains outside the formal apparatus, consistent with
Remark~\ref{rem:not_replacing}.

\subsubsection{Becker 1965: household production, constrained}

Becker's \emph{A Theory of the Allocation of
Time}~\citep{becker1965theory} introduces household production
functions $z_j = F_j(x_{j,1}, \dots, x_{j,n_j})$ that produce
utility-yielding commodities from market-good and time inputs.
The production functions are unconstrained: any specification of
$F_j$ is admissible, and the framework's empirical content depends
on what specification the modeller chooses.

GFM's M6 (cooperative capabilities) is the formal counterpart to
Becker's household production, with one crucial constraint: the
cooperative excess admitted by M6 is poset-coherent, bounded by
the structural connectivity of the capability poset rather than
by an arbitrary $F_j$.  This is what converts Becker's
descriptively-flexible mathematics into a structurally-falsifiable
framework.  Predictions~\ref{pred:bundle_completion}
and~\ref{pred:coop_premium} depend on the constraint: the
poset-coherence of cooperative excess is what makes the
predicted variance pattern in Prediction~\ref{pred:coop_premium}
testable.  Channel~4 (bundle decomposition) is the operational
expression of the constraint at the per-event level.

\subsubsection{Bernheim--Rangel 2009: behavioral welfare economics' structural neighbor}

\citet{bernheim2009beyond} develop behavioral welfare economics
on a choice-theoretic foundation that distinguishes
welfare-relevant from welfare-irrelevant choice domains: when
agents exhibit dynamically inconsistent or context-dependent
choices, the framework asks which subset of the agent's
revealed-choice data carries welfare-relevant information.  The
construction is a structural neighbor of GFM's polarity
boundary: both frameworks separate revealed-sacrifice events
into two operational regimes: one in which the trade reveals
welfare-relevant content, one in which it does not.  The
divisions are different in kind: Bernheim--Rangel's split is
across choice domains (the choice was made under conditions
that compromise welfare-revelation), while GFM's polarity
boundary is across capability sufficiency (the trade is on a
need-side or want-side dimension), but the operational
discipline (``not all revealed sacrifice reveals
preference'') is shared.  The closest historical parallel for
the welfare-revelation distinction is Kahneman's
experienced-utility-vs-decision-utility
distinction~\citep{kahneman1997back}: a domain-of-validity
restriction on revealed-preference inference.

GFM does not adopt the Bernheim--Rangel apparatus directly; the
polarity boundary of \citep[Need~Sufficiency]{lasser2026paper8b}
operates on capability sufficiency rather than on choice-domain
welfare-revelation.  But the structural overlap is real: both
frameworks recognize that revealed sacrifice is not always a
preference signal, and both supply an operational discipline
for separating preference-revealing from preference-obscuring
trades.  A complete formal comparison (whether the polarity
boundary plus the S1-attestation layer can be derived as a
specialization of Bernheim--Rangel's choice-domain restriction,
or whether the two are structurally distinct) is open and out
of scope for this paper.

\subsubsection{Stiglitz--Sen--Fitoussi 2009: the gap decomposition, as a measurement instrument}

The 2009 SSF report~\citep{stiglitz2009report} is the most
direct qualitative precursor.  Its central recommendation
(``GDP measures poorly capture welfare because they miss
non-market production, quality-of-life capabilities, and
distributional equity'') is a qualitative gap-decomposition
statement: the
welfare-relevant capability space exceeds the GDP-measured
subspace, and the difference admits structurally distinct
components (non-market production, quality of life, distribution).
The report's recommended remedies (time-use surveys, satellite
accounts, multi-dimensional well-being indices) are imputation
moves that re-individuate previously-uncounted activities into the
measured space.

\citep[Need~Sufficiency]{lasser2026paper8b}'s
five-cell $\delta$-decomposition (Definition~8) is the
quantitative formalization.  $\delta_{\mathrm{covered}}$,
$\delta_{\mathrm{dormant}}$, $\delta_{\mathrm{restricted}}$,
$\delta_{\mathrm{residual}}$, and $\delta_{\mathrm{boundary}}$
are the structural cells SSF was identifying qualitatively.
$\ResS$ is the formal version of ``the parts GDP misses'';
re-individuation (Channel~3) is the formal version of imputation
methods.

A measurement-instrument observation worth making here.  In
welfare-economics practice, the SSF decomposition is itself a
\emph{measurement instrument}: it tells the analyst which
imputation method applies to which kind of unobserved activity
(time-use surveys for household labour, satellite accounts for
ecosystem services, imputation for owner-occupied housing).  The
decomposition is not a separate measurement primitive that adds
to the ledger; it is a deployment-readable diagnostic that
selects which intervention to apply.  In the GFM apparatus, the
five-cell $\delta$-decomposition plays the same selector role:
when a deployment reports a wide $\epsgap$, the cell-wise
attribution tells the deployer whether the binding channel is
observation density (large $\delta_{\mathrm{dormant}}$),
attestation (large attestation-limited share of
$\delta_{\mathrm{covered}}$), individuation (large
$\delta_{\mathrm{residual}}$), or bundle decomposition (large
intra-bundle mismatch share of $\delta_{\mathrm{covered}}$).
The lineage from SSF's qualitative decomposition to GFM's
quantified decomposition is therefore not just structural but
operational: both objects play the same epistemic role in their
respective measurement traditions, and the GFM version inherits
the role formally.

\subsubsection{Alkire--Foster 2011: the closest existing quantitative apparatus}

The capability-welfare-economics tradition has produced one
substantially-quantitative measurement framework that the
preceding subsections did not engage: the
\emph{Alkire--Foster} multidimensional-poverty
methodology~\citep{alkire2011counting}.  Alkire--Foster
operates on a vector of capability dimensions (each with a
deprivation threshold), counts deprivations per agent against
those thresholds, and aggregates into a multidimensional
headcount or intensity-weighted index.  This is the quantitative
apparatus closest in operational structure to GFM's combination
of capability poset, polarity boundary, and gap decomposition.

Two structural differences distinguish GFM from the
Alkire--Foster framework, both of which the lineage claim
should make explicit rather than gloss.  First, Alkire--Foster
aggregates by counting (or weighted counting) of dimension-wise
deprivations; GFM aggregates by a poset measure $\volL$ that
admits cooperative excess (M6) and is not in general reducible
to a sum or weighted sum of dimension-wise contributions.  The
two aggregations coincide only when cooperative capabilities
vanish, a fungibility-collapse-adjacent regime in which
multidimensional deprivation reduces to additive deprivation,
which Alkire--Foster does not commit to but in practice
operates within.  Second, Alkire--Foster's deprivation
thresholds are normatively chosen (per the application) and the
framework does not formalize the structural distinction between
``cost-of-access below threshold'' and ``preference above
threshold''; \citep[Need~Sufficiency]{lasser2026paper8b}'s
polarity boundary is the GFM-side structural addition that
Alkire--Foster's apparatus could in principle inherit.

The relationship is therefore: Alkire--Foster is a
quantitative welfare-measurement framework that GFM contains
under additional structural commitments (cooperative capabilities,
polarity discipline, residual-class characterization).  The
positioning of Alkire--Foster relative to GFM is closer to that
of standard utility theory under fungibility collapse
(Theorem~\ref{thm:fungibility_collapse}) than to the qualitative
tradition of the prior subsections: the Alkire--Foster framework
is already quantitative, and GFM's contribution is the
additional structural axiomatization rather than a
qualitative-to-quantitative formalization.  A complete formal
comparison, including whether a ``cooperative-coherent
Alkire--Foster index'' would be a useful intermediate
methodology, is open; we flag the lineage but do not develop it
further here.

\subsubsection{Basu--L\'opez-Calva 2011: prior formalization-survey of capability sets}

\citet{basu2011functionings}, in the \emph{Handbook of Social
Choice and Welfare}, survey prior formal-mathematical
treatments of capability sets, covering a wide range of
approaches including Sen/Foster axiomatizations, Herrero/Roemer
cardinality and indirect-utility constructions, fuzzy-set
aggregations, Suzumura-Xu extended-alternatives frameworks, and
operationalizations such as the Human Development Index.  This
is the canonical reference for the prior formal-mathematical
landscape on the capability approach, and its placement in the
Handbook of Social Choice and Welfare establishes the territory
GFM operates within.

The Basu--L\'opez-Calva survey covers approaches that share with
GFM a commitment to formal apparatus on the capability set, but
diverge from GFM in the structural objects used.  Three
representative branches illustrate the structural distance:
\begin{itemize}
\item \emph{Cardinality-based formalizations} rank capability
  sets by counting elements (with various refinements for
  weighting).  GFM's poset-with-measure approach reduces to a
  weighted cardinality only when M6's cooperative excess
  vanishes; in general the volume measure is not a
  cardinality.
\item \emph{Indirect-utility ranking} compares capability sets
  by the maximal utility achievable on each.  This collapses
  the capability set to a single utility number, which GFM
  explicitly avoids: $\volL$ is a measure on the set, not a
  best-element selector.
\item \emph{Effective-freedom and extended-alternative
  relations} rank capability sets by various non-cardinal
  axiomatic conditions.  These are closer in spirit to GFM
  than the first two branches but remain non-measure-theoretic:
  they produce orderings, not measures.
\end{itemize}
GFM's measure-theoretic-poset framework is structurally
distinct from all three approaches Basu--L\'opez-Calva survey:
the apparatus produces a measure (not a ranking) on a poset
(not a flat set or chain), with cooperative excess (M6) and a
polarity boundary as first-class structure that the surveyed
approaches do not formalize.  The closest comparator within the
surveyed literature is the effective-freedom relations branch;
GFM's contribution there is to upgrade the ordering to a
measure with explicit axiomatic structure.

\subsubsection{Pattanaik--Xu / Sugden: opportunity-set axiomatic neighbors}

A related literature is the axiomatic-opportunity-set tradition,
which ranks opportunity sets by axiomatic conditions on the
ranking relation~\citep{pattanaik2000ranking}.  This literature
is structurally adjacent to the capability approach (capability
sets are formally a kind of opportunity set), and Sugden's
``Opportunity as a Space for Individuality: Its Value and the
Impossibility of Measuring It''~\citep{sugden2003opportunity}
proves an impossibility result on measuring
opportunity-as-individuality directly: under reasonable axioms,
no real-valued measure on opportunity sets can simultaneously
satisfy the axioms.  This is a prima facie objection to a
volume-measure-on-capability-poset framework, and the lineage
should engage with it.

The objection does not automatically apply to GFM, for a
structural reason worth naming.  Sugden's impossibility is on
measures of opportunity-as-individuality, i.e., the
freedom-to-choose dimension treated as the welfare-relevant
quantity.  GFM's
volume measure is on \emph{capabilities exercised}
($\volRwin$) and \emph{capabilities possessed} ($\volL$), with
the freedom-to-choose dimension recovered through the polarity
boundary (\citep[Definition~3]{lasser2026paper8b}) and the
S1-attestation layer (\citep[Remark~18]{lasser2026paper8a})
rather than through a single welfare-relevant
opportunity-as-individuality scalar.  GFM measures realized
exercise plus capability volume; Sugden's impossibility applies
to opportunity-as-individuality.  The two notions are
structurally distinct on the exercise side, but $\volL$ remains
a scalar measure on possessed-capability volume, a slice on
which Sugden-style concerns may still apply, and the
sidestep claim is therefore conditional rather than complete.
A complete proof that GFM sidesteps the relevant Sugden-style
impossibility result (rather than falling under it on the
$\volL$-slice or on a different formal slice) is the extension
of Open~Question~\ref{oq:arrow} to the opportunity-set
literature.  The qualitative argument is sound on the exercise
side; the formal version, including the $\volL$-slice, is
open.

\subsubsection{Arrow 1951: the impossibility GFM appears to sidestep}

Arrow's impossibility theorem~\citep{arrow1951social}
establishes that no social welfare function aggregating
individual ordinal preferences into a social ordering can
simultaneously satisfy all five Arrovian axioms (transitivity,
unrestricted domain, Pareto, independence of irrelevant
alternatives, and non-dictatorship).  Subsequent work
(\citep{sen1970collective} and a substantial literature) explores
ways out: relaxing IIA, adopting cardinal utility, or restricting
the domain.  \citet{adler2012wellbeing} develops an axiomatized
prioritarian social welfare function as a modern formal-aggregation
route around Arrow, a live alternative to the cardinal-utility
move of \citep{sen1970collective} and the Sugden-style
opportunity-set axiomatics
of~\citep{sugden2003opportunity}.

GFM \emph{appears} to sidestep the impossibility by being a
different kind of aggregation.  $\volL$ is not a ranking
function; it is a measure on a poset.  It does not aggregate
ordinal preferences into a social ordering; it aggregates
capabilities into a capability volume.  The Arrovian axioms have
only partial analogs:
\begin{itemize}
\item \textbf{Transitivity:} satisfied trivially (real-valued
  measures are totally ordered on their values).
\item \textbf{Unrestricted domain:} $\volL$ is defined on any
  poset satisfying~M1--M6
  (\citep[Proposition~1]{lasser2026poset}).
\item \textbf{Pareto:} the monotonicity component of M3 ensures
  that strictly increasing capability strictly increases $\volL$.
\item \textbf{Independence of irrelevant alternatives:} does not
  apply to a volume measure.  IIA is defined via pairwise
  comparisons in a ranking, and $\volL$'s value at a capability
  set does not depend on pairwise comparisons with alternative
  sets.  The axiom is not satisfied or violated; it is
  inapplicable to the formulation.
\item \textbf{Non-dictatorship:} $\volL$ aggregates over agents
  via the shared poset; no single agent's preferences determine
  the aggregate.  The individuation discipline of M5 is what
  makes this an aggregation rather than a comparison; the
  per-agent benchmark calibrations $w_i(\cdot)$ on the shared
  poset (Remark~\ref{rem:heterogeneity}) are the formal
  expression.
\end{itemize}
GFM's apparent avoidance of Arrow's original theorem is
therefore structural, not axiomatic: the framework does not
perform the kind of aggregation Arrow's original theorem
applies to.  This is an alternative way out from the
cardinal-utility move of \citep{sen1970collective} and the
prioritarian-SWF axiomatization
of~\citep{adler2012wellbeing}; all three are forms of relaxing
the ordinal-preference-aggregation framing, with each making a
different structural commitment.  Sen's relaxation adds
cardinal information to the preferences; Adler's adds priority
weights to the social welfare function; GFM's switches to a
measure-theoretic apparatus on a poset that does not aggregate
preferences at all.  Whether the relaxed setting admits its own
impossibility analogue remains open; see
Open~Question~\ref{oq:arrow}.

\paragraph{What this paper does \emph{not} establish.}
Sidestepping a specific impossibility result is not the same as
proving no analogous impossibility holds.  Whether the
measure-on-poset formulation admits its own impossibility-style
result (some incompatible combination of axioms on volume
aggregation that the framework's M1--M6 silently violate, or an
unrestricted-domain analog that produces a contradiction with
the cooperative-superadditivity axiom~M6) is not addressed in
this paper.  We have shown that the Arrovian formulation does
not directly apply (IIA is inapplicable, the other axioms have
partial analogs that GFM satisfies), not that no Arrow-style
result applies in the new formal setting.  This is
Open~Question~\ref{oq:arrow}, and a complete treatment requires
either a no-impossibility proof on the measure-on-poset class
or an honest acknowledgment of an analogous limitation.  Sen's
own work has been congenial to this kind of move (the
capability approach was already a step away from the
preference-aggregation framing), but the formal completion is
open.

\subsubsection{Summary mapping}

Table~\ref{tab:lineage_mapping} consolidates the lineage across
three categories.  The qualitative-tradition precursors (Sen,
Nussbaum, Becker, SSF) supply moves the channels formalize.  The
already-quantitative comparator (Alkire--Foster) supplies the
nearest existing quantitative apparatus that GFM contains under
additional structural commitments.  The structural neighbors
(Bernheim--Rangel on behavioral welfare economics,
Basu--L\'opez-Calva on prior capability-set formalizations,
Sugden / Pattanaik--Xu on opportunity-set axiomatics, Arrow /
Adler on impossibility and prioritarian SWF) situate GFM
within the broader formal-welfare-economics landscape rather
than mapping cleanly onto specific channels.

\begin{table}[ht]
\centering
\footnotesize
\setlength{\tabcolsep}{4pt}
\begin{tabular}{>{\raggedright\arraybackslash}p{0.13\linewidth}
                >{\raggedright\arraybackslash}p{0.30\linewidth}
                >{\raggedright\arraybackslash}p{0.32\linewidth}
                >{\raggedright\arraybackslash}p{0.16\linewidth}}
\toprule
\textbf{Precursor} & \textbf{Qualitative content} &
\textbf{GFM quantification} & \textbf{Channel} \\
\midrule
Sen 1985 & Capability set, functioning vector, capability/functioning gap &
Poset $P$, $\volL$, $\volRwin$, five-cell $\delta$-decomposition &
All four \\
\addlinespace
Nussbaum 2000 & Basic vs higher capabilities (normative) &
Polarity boundary + downstream-safety condition &
Channel~2 \\
\addlinespace
Becker 1965 & Household production $F_j$ (unconstrained) &
M6 cooperative caps (poset-coherent) + $\alpha_{n,c}$ &
Channel~4 \\
\addlinespace
Bernheim--Rangel 2009 & Choice-domain welfare-revelation distinction (behavioral) &
Polarity boundary + S1-attestation as capability-sufficiency split &
Channel~2 \\
\addlinespace
SSF 2009 & GDP-misses-non-market gap (qualitative) &
$\ResS$ + five-cell $\delta$-decomp (deployment-readable) &
Channels~2, 3 \\
\addlinespace
Alkire--Foster 2011 & Multidimensional poverty index (counting / weighted-sum aggregation) &
M6 cooperative excess + polarity boundary as additional structure &
n/a (already quantitative) \\
\addlinespace
Basu--L\'opez-Calva 2011 & Survey of capability-set formalizations (cardinality, indirect-utility, effective-freedom) &
Poset measure with M6 + polarity as structural upgrade beyond surveyed approaches &
n/a (formalization survey) \\
\addlinespace
Sugden 2003 / Pattanaik--Xu 2000 & Opportunity-set ranking axiomatics; Sugden's measurement-impossibility &
GFM measures realized exercise + capability volume, not opportunity-as-individuality &
n/a (structural neighbor) \\
\addlinespace
Arrow 1951 / Adler 2012 & Ordinal-preference impossibility; prioritarian SWF &
$\volL$ as measure on poset (appears to sidestep original framing); not a preference aggregator &
n/a \\
\bottomrule
\end{tabular}
\caption{The capability-economics lineage mapped onto GFM
quantifications and the four channels of
Sections~\ref{sec:channel_observation}--\ref{sec:channel_bundles}.
The SSF~2009 row also flags the diagnostic-decomposition role:
the five-cell $\delta$-attribution acts as a deployment-readable
selector for which channel binds, mirroring the role the SSF
decomposition plays in welfare-economics measurement
methodology.  The Arrow~1951 row is structural rather than
channel-bound: GFM appears to sidestep the impossibility by
being a measure on a poset rather than a ranking aggregator,
but whether an analogous impossibility result applies to
measure-aggregation on posets is open
(Open~Question~\ref{oq:arrow}).}
\label{tab:lineage_mapping}
\end{table}

The historical reading the lineage supports: GFM is not a new
welfare-economics framework competing with Sen, Nussbaum, Becker,
or SSF.  It is a measure-theoretic formalization of the same
tradition, and the alignment-tightening channels are exactly the
operational moves the qualitative tradition has been negotiating.
The novel contribution this paper adds is the connection between
the formalization and the alignment programme: the same channels
the welfare-economics tradition uses to close measurement gaps are
the channels that, under the Goodhart machinery of
Section~\ref{sec:goodhart}, tighten alignment guarantees in the
direction of the operational truth.  Better economic models are
not just better welfare measurement; they are tighter alignment
guarantees, and the tightening mechanism is the same one welfare
economics has been refining since SSF 2009.
